\section{Effective algebraic-circle bounds}\label{sec:arithmetic53}
For algebraic unit-circle eigenvalues of infinite order, an explicit separation estimate replaces the Diophantine hypothesis of the preceding section. The parameters $r,m,A_r$ below are those of \eqref{eq:r53}--\eqref{eq:circleconstants53}; in this section $r$ denotes the localization order, not the number of command rotations. The alphabet contains $I,R_\theta$, and one interface mean is $\rho\cos\phi$.

\begin{lemma}[Separated powers give a block defect]\label{lem:separated53}
Suppose $|\lambda^n-1|\ge\sigma$ for every integer $1\le n\le2mk$, where $0<\sigma\le1$. Write $R_m(t)=\diag(R_t,R_{2t},\ldots,R_{mt})$ for the representation acting on $Z_m$. On this representation, the uniform law on the $2k$ words with products $R_\theta^j$, $0\le j<2k$, each padded to length $2k$, has defect at least $\sigma^2/16$ in \eqref{eq:gap54}.
\end{lemma}
\begin{proof}
Write a unit vector as $u=(u_1,\ldots,u_m)$ in its two-dimensional frequency blocks. For distinct indices $i,j\in\{0,\ldots,2k-1\}$,
\[
 \norm{R_m(i\theta)u-R_m(j\theta)u}^2
 =\sum_{l=1}^m\norm{u_l}^2|\lambda^{l(i-j)}-1|^2\ge\sigma^2.
\]
A ball of radius strictly less than $\sigma/2$ about any of $k$ centers contains at most one of these $2k$ points. Thus at least $k$ points have distance at least $\sigma/2$ from all centers (the same conclusion follows by taking the open balls of radius $\sigma/2$). For unit centers the defect is one half of the squared distance to the closest center. Averaging gives at least $(1/2)(\sigma^2/8)=\sigma^2/16$.
\end{proof}

\begin{theorem}[Full-noise algebraic-circle bound]\label{thm:algebraic53}
Assume $\lambda$ is algebraic and not a root of unity. Let its primitive irreducible integer polynomial be
\[
 P(z)=a z^d+a_{d-1}z^{d-1}+\cdots+a_0,\qquad a>0.
\]
Set $H_P=\max(a,|a_{d-1}|,\ldots,|a_0|)$ and
\[
                     B_P=a(1+H_P/a)^{d-1}>1.
\]
With $r,m,A_r$ as in \eqref{eq:r53}--\eqref{eq:circleconstants53}, every error-$\varepsilon$ realization of peak at most $k$ satisfies
\begin{equation}\label{eq:algwidth53}
 N<2k\left(1+16\,2^{2(d-1)}B_P^{4mk}\log A_r\right).
\end{equation}
Its number of command cuts $t\in\{1,\ldots,N\}$ with width at most $k$, even if other widths are unrestricted, is at most
\begin{equation}\label{eq:algoccupation53}
 2k-1+32k\,2^{2(d-1)}B_P^{4mk}\log A_r.
\end{equation}
Consequently, for fixed algebraic data, $\rho$, and $\varepsilon<\rho/2$,
\begin{equation}\label{eq:algrate53}
 W_{N,\varepsilon}\ge
 \frac{\log N-O(\log\log N)}{4m\log B_P}.
\end{equation}
All constants in the finite inequalities are explicit. If $\lambda=(a_1+ib_1)/c$ with integers $a_1,b_1$, $c>0$, $a_1^2+b_1^2=c^2$, and $\lambda$ not a root of unity, then \eqref{eq:algwidth53}--\eqref{eq:algrate53} hold with the factor $2^{2(d-1)}$ deleted and $B_P$ replaced by $c$.
\end{theorem}
\begin{proof}
List the conjugates with $\lambda_1=\lambda$. They are distinct by irreducibility in characteristic zero. None is a root of unity: otherwise irreducibility would make $P$ a cyclotomic polynomial up to its primitive sign. The resultant
\[
 \operatorname{Res}(P,z^n-1)=a^n\prod_{j=1}^d(\lambda_j^n-1)
\]
is therefore a nonzero integer. Each conjugate has modulus at most $R_P=1+H_P/a$. For completeness, if $|z|>R_P$, then
\[
 \sum_{j=0}^{d-1}|a_j||z|^j
 <\frac{H_P|z|^d}{|z|-1}<a|z|^d,
\]
so $P(z)\ne0$. Using $|z^n-1|\le2\max(1,|z|)^n$ for each other conjugate, the resultant bound gives
\begin{equation}\label{eq:resultant53}
 |\lambda^n-1|\ge
 2^{-(d-1)}\left(a\prod_{j=2}^d\max(1,|\lambda_j|)\right)^{-n}
 \ge 2^{-(d-1)}B_P^{-n}.
\end{equation}
For $1\le n\le2mk$, take $\sigma=2^{-(d-1)}B_P^{-2mk}$ in Lemma~\ref{lem:separated53}. The block defect is at least
\[
                  d_k=\frac{2^{-2(d-1)}B_P^{-4mk}}{16}.
\]
Repeated conditional contraction and Lemma~\ref{lem:harmonic53} imply
\[
 \left\lfloor\frac{N}{2k}\right\rfloor
 \le \frac{\log A_r}{-\log(1-d_k)}
 \le 16\,2^{2(d-1)}B_P^{4mk}\log A_r.
\]
This proves \eqref{eq:algwidth53}. The greedy occupation argument in the proof of Theorem~\ref{thm:radius54}, with $b=0$ and block length $2k$, proves \eqref{eq:algoccupation53}. Taking logarithms gives \eqref{eq:algrate53}: fix $C>1/(4m\log B_P)$. If $k\ge C\log N$ the claim is immediate; otherwise $\log k\le\log C+\log\log N$, which can be absorbed in the logarithm of \eqref{eq:algwidth53}. The implicit constant may depend on $P,\rho,\varepsilon$, not on $N$.

In the rational-circle case, $(a_1+ib_1)^n-c^n$ is a nonzero Gaussian integer, of modulus at least one. Thus $|\lambda^n-1|\ge c^{-n}$ directly. Repeat the proof with $\sigma=c^{-2mk}$. Here $c>1$ automatically. A nonminimal supplied denominator is still valid but weakens the bound.
\end{proof}
The threshold is sharp; these width rates are not asserted sharp. The factor $m$ records the increased localization degree as the error approaches the boundary. At smaller errors, the retained coordinate-calibration estimates can have better constants. One should use the larger of applicable lower bounds, rather than claiming numerical improvement at every parameter value.

\begin{proposition}[An explicit nonrational algebraic example]\label{prop:salem53}
Let $t=(1-\sqrt{13})/2$ and
\[
 \lambda=\frac{t+i\sqrt{4-t^2}}2.
\]
Then $|\lambda|=1$, $\lambda$ is not a root of unity, and one can take $d=4$ and $B_P=8$ in Theorem~\ref{thm:algebraic53}.
\end{proposition}
\begin{proof}
The number $t$ lies strictly between $-2$ and $2$ and solves $t^2-t-3=0$. Substitution of $t=z+z^{-1}$ gives
\[
                       P(z)=z^4-z^3-z^2-z+1.
\]
Its reduction modulo two is $z^4+z^3+z^2+z+1$. It has no linear factor over $\mathbb F_2$ and is not divisible by the only monic irreducible quadratic $z^2+z+1$; hence it is irreducible. Gauss's lemma proves irreducibility over $\mathbb Q$. The conjugate value $(1+\sqrt{13})/2>2$ gives a positive real conjugate of $\lambda$ greater than one. A root of unity cannot have such a conjugate. Finally $a=H_P=1$ and $d=4$, so $B_P=2^3=8$.
\end{proof}
For a reproducible rational calibration, take $\rho=1/10$ and $\varepsilon=1/25$. Then $\delta=2/25$, $r=5$, $m=6$, $\Delta_r=1/150$, and $A_r=23364$. Thus the rational-circle alphabet above obeys
\[
                 N<2k\bigl(1+16\cdot5^{24k}\log23364\bigr).
\]
This is a concrete bound inside the interval between the former small-error obstruction and the sharp threshold $1/20$. The exact coefficient and resultant checks accompanying the article verify the displayed arithmetic, not the universal theorem by finite experimentation.

